% Generated draft: requires mathematical and source review.
\documentclass[11pt]{article}
\usepackage[margin=25mm]{geometry}
\usepackage{fontspec}
\setmainfont{Latin Modern Roman}
\usepackage{amsmath,amssymb,mathtools,unicode-math}
\setmathfont{Latin Modern Math}
\usepackage{enumitem,needspace,xurl,hyperref}
\setlength{\parindent}{0pt}
\setlength{\parskip}{5pt}
\setlength{\emergencystretch}{3em}
\newcommand{\N}{\mathbb{N}}
\newcommand{\Z}{\mathbb{Z}}
\newcommand{\Q}{\mathbb{Q}}
\newcommand{\R}{\mathbb{R}}
\newcommand{\C}{\mathbb{C}}
\newcommand{\F}{\mathbb{F}}
\newcommand{\A}{\mathbb{A}}
\newcommand{\PP}{\mathbb{P}}
\newcommand{\E}{\mathbb{E}}
\newcommand{\eps}{\varepsilon}
\newcommand{\dd}{\,\mathrm{d}}
\begin{document}
\section*{UnsolvedMath ID: 30004589}
\textbf{Source problem code:} catalogue entry
\subsection{Definitions and mathematical statement}
For a positive integer \(n\), let \(\mathbb Z_n=\Z/n\Z\) and let \(T_n=\mathbb Z_n^2\) be the two-dimensional discrete torus, with periodic boundary conditions. An Ising spin configuration on \(T_n\) is an element \(\sigma\in\{-1,+1\}^{T_n}\).

\textbf{Problem as supplied.} Determine the exact reconstruction threshold for the critical Ising model on \(T_n\).

\textbf{Unresolved formulation.} The supplied record does not define reconstruction: it specifies neither the observations available, the object to be reconstructed, the admissible reconstruction procedures, nor the success criterion. It also does not identify the parameter whose threshold is sought or specify its limiting regime as \(n\) varies. The interaction, external field, and meaning of criticality are not explicitly stated. Consequently, the record does not determine a unique quantified threshold problem. These specifications must be recovered from the intended source before a precise reconstruction statement can be given; no particular reconstruction scheme or additional model hypotheses are imposed here.
\subsection{Short English statement}
The question asks for an exact reconstruction threshold for the critical Ising model on a two-dimensional torus. The supplied record leaves the observation scheme, reconstruction target, success criterion, and threshold parameter undefined, so a precise threshold formulation requires further source clarification.
\subsection{Sources}
\begin{itemize}
\item Local UnsolvedMath registry, permanent ID \(30004589\), source problem code \texttt{catalogue entry}, title ``Reconstruction Threshold for Critical Ising Models on the Torus.'' The supplied record lists proposed year 2021 and names no proposer.
\item Supplied report citation: \emph{Spatial Networks and Percolation} (2022), \emph{Oberwolfach Reports}, pp.\ 149--224, DOI 10.4171/owr/2021/4. \url{https://doi.org/10.4171/owr/2021/4}. The background describes these reports as documenting mathematical workshops and their open-problem discussions.
\item Supplied literature reference: Galicza and Pete, \emph{Ising and factor measures}. \url{https://arxiv.org/abs/2406.09232}.
\item Literature assessment dated exactly \texttt{2026-08-22}: the supplied curation labels the problem open and states that subsequent work gives nontrivial upper and lower reconstruction regimes for critical two-dimensional Ising boxes, while the exact threshold requested here remains unresolved. This is supplied dated curation, not an independently verified literature assessment. Its time-sensitive status and cited claims require independent verification; the reference to boxes does not itself specify a torus formulation.
\item The registry labels the statement \texttt{corrected\_verified}; its supplied verification explanation concerns dollar-delimiting explicit formulas and preservation of the supplied mathematical content. This does not establish completeness or mathematical correctness of the reconstruction formulation.
\end{itemize}
\end{document}
